\section{Prospective published-score evaluation}
\label{sec:new-evaluation}
\subsection{Inputs, fixed comparisons, and independent controls}
The MSAD archive supplies twelve stored detector configurations:
IFOREST, LOF, MP, NORMA, IFOREST1, HBOS, OCSVM, PCA, AE, CNN, LSTM and
POLY~\citep{msad2023}. Before computing metric outcomes, we fixed a
family-stratified SHA-256 order over 1,863 metadata-eligible files,
excluding training filenames and raw members larger than 12 MB. A fixed
round-robin queue and validity/identity rules selected 128 files from
sixteen families. Of 148 considered files, nineteen identity duplicates
and one single-class file were excluded; no acquisition failure changed
the selection. The final panel has 15,793,514 temporal observations and
1,536 score vectors. A separate parser rechecked all 1,664 source members
and 1,792 arrays, including 189,522,168 score scalars.

All arrays are preserved without padding, repair, new quantization or
outcome-based exclusion. Each stored score vector has at most 101 distinct
values; every textual score has at most two decimal places. Thus this is
a study of released, quantized scores, not recovered full-precision model
outputs. The budget $H=250$ already exceeds the number of distinct scores,
but uniform \emph{rank} sampling can revisit large ties and miss rare score
classes. Enumerating every distinct score is a simple, important control.

The primary setting is $L=30,H=250$; widths 100 and 300 and budgets
32, 64, 128, 500 and 1,000 are sensitivity analyses. We retain full tied
winner sets at tolerance $10^{-10}$, and report best- and worst-case
exhaustive regret among sampled tied winners. The prospective practical
decision uses
\[
 \widehat{\mathcal W}=\{m:s_m\geq\max_j s_j-10^{-10}\},\qquad
 \Delta_{\rm best}=\max_j e_j-\max_{m\in\widehat{\mathcal W}}e_m,
\]
where $e_m$ and $s_m$ are the exhaustive and sampled values. Worst-case
regret replaces the second maximum by a minimum. The prospective practical
decision rule requires best-case regret at least 0.01 in at least 5\%
of series and at least three source families. This rule is a declared
finite-panel study threshold, not a universal utility cutoff or publication
criterion. Per-family and equal-family results accompany pooled results;
model pairs and widths are not treated as independent observations.

At the primary setting, all 1,536 matched sampled results agree with the
unchanged pinned upstream implementation to $1.06\times10^{-15}$.
The fixed first two retained series per family give 30 validation series
and 360 score vectors. All agree with both an independent dense exhaustive
implementation and unchanged upstream evaluation at $H=n$, with maximum
absolute differences $4.45\times10^{-15}$ and $3.34\times10^{-16}$.
All 720 independent exhaustive calls complete without errors or timeouts.
These are numerical controls; integer evidence below strengthens the
ordering and regret claims to the ideal real-arithmetic target.

\subsection{Sampling changes model-selection decisions}
At $L=30,H=250$, the sampled winner set changes on 10/128 series, including
nine changes between unique winners. Eight series (6.25\%) in four
families have best-case exhaustive regret at least 0.01; the maximum
regret is 0.398575. Thus the prespecified finite-panel effect threshold
is met. The equal-family practical-error fraction is 5.28\%.
A 2,000-replicate source-family bootstrap gives a pooled 95\% percentile
interval of [1.09\%, 13.20\%]. The observed panel passes the declared
threshold; this interval does not establish population prevalence above
5\%. Table~\ref{tab:primary-families} retains every family, including the
twelve with no practical winner error.

\begin{table}[t]
\centering\small
\begin{tabular}{lrrr@{\qquad}lrrr}
\toprule
Family & $N$ & errors & max regret & Family & $N$ & errors & max regret\\
\midrule
Daphnet & 7 & 0 & .0000 & MITDB & 8 & 0 & .0000\\
Dodgers & 1 & 0 & .0000 & NAB & 9 & 0 & .0000\\
ECG & 11 & 0 & .0000 & OPPORTUNITY & 10 & 0 & .0000\\
GHL & 10 & 1 & .0272 & Occupancy & 2 & 0 & .0000\\
Genesis & 1 & 0 & .0000 & SMD & 9 & 0 & .0000\\
IOPS & 11 & 0 & .0000 & SVDB & 10 & 0 & .0000\\
KDD21 & 10 & 2 & .3986 & SensorScope & 10 & 0 & .0000\\
MGAB & 10 & 1 & .2517 & YAHOO & 9 & 4 & .2362\\
\bottomrule
\end{tabular}
\caption{Primary setting on all 128 series. ``Errors'' counts sampled
winner sets whose \emph{best} exhaustive value is at least 0.01 below the
exhaustive maximum. Maximum regret includes all cases in each family.}
\label{tab:primary-families}
\end{table}

Errors have both signs. Across the 1,536 scalar evaluations, mean absolute
error is 0.009286, maximum underestimation is 0.560123 and maximum
overestimation is 0.161692. Fifteen values overestimate by at least 0.01
and 150 underestimate by at least 0.01. There are 106 strict pairwise
reversals, 22 with both exact and sampled margins at least 0.01; these
pairs are descriptive, not independent samples. The new inflation mechanism
is therefore not a universal explanation for the observed selection errors.

\paragraph{Source-overlap sensitivity.}
The selection excluded exact label, observation and standardized-observation
identities against 362 prior inputs. A subsequent bounded audit additionally
compared joint observations and labels after dropping either endpoint. It
found 20/128 matches to prior inputs under these transformations. We keep
the frozen 128-case primary panel and report the descriptive 108-case
nonmatching subgroup: four cases (3.70\%) in two families have regret at
least 0.01, with maximum 0.236208. This subgroup does not meet the original
5\%/three-family threshold. The 20 matched cases contain four practical
errors and the overall maximum. The bounded audit neither proves the
remaining physical recordings independent nor permits describing the
full panel as previously unseen recordings. Its newly evaluated content
is the released detector-score bank under this pinned metric contract.

\subsection{A measured instance of post-saturation inflation}
To connect the construction to actual scores, an exploratory diagnostic
uses all twelve models on all eight primary harmful-selection series.
For each width it inserts the first exhaustive threshold at which recall
reaches one into the original sampled grid, adding at most one threshold.
Integer mass brackets independently certify each first-saturation event
and its predecessor on all 96 score vectors at all 31 widths.
The difference between the original and inserted-grid areas measures the
net effect of postponing that saturation boundary; the remaining
inserted-grid error measures other omitted-threshold effects. This is an
outcome-selected mechanism analysis, not a second prevalence estimate or
a claimed faster algorithm.

The largest-regret series, KDD21/005 (the distorted CIMIS air-temperature
recording), has 8,183 points and one original positive. Its sampled winner
is LOF, while exhaustive evaluation selects CNN
(Table~\ref{tab:real-mechanism}). For LOF, inserting saturation reduces the
area by 0.129093, or 79.84\% of its total 0.161692 overestimation.
The net change decomposes into a positive 0.181891 buffer-mass term,
zero original-mass term and a negative 0.052797 count-dilution term.
For CNN, count dilution dominates: the corresponding terms are
$+0.026004$ and $-0.470588$, explaining its 0.444584 underestimation.
Inserting saturation recovers the exhaustive winner on three of the eight
series; five selections remain wrong. Thus the mechanism occurs in real
published scores, but correcting it alone is insufficient in general.

\begin{table}[t]
\centering\small
\begin{tabular}{lrrr}
\toprule
Model & Exhaustive & Sampled, $H=250$ & Saturation inserted\\
\midrule
CNN & .500000 & .055416 & .500000\\
LOF & .101425 & .263117 & .134024\\
\bottomrule
\end{tabular}
\caption{Two unchanged published score vectors on KDD21/005. The diagnostic
inserts at most one threshold per width and isolates a particular error
component; it does not remove all integration error.}
\label{tab:real-mechanism}
\end{table}

\subsection{Integer certificates and a full-grid control}
All 3,072 partial-grid certificate calls complete: twelve configurations
on all 128 series at $L=30$, for $H=250$ and $H=1,000$. A model is certified
as the unique exhaustive winner only if its integer lower endpoint exceeds
every competitor's integer upper endpoint. At $H=250$, 63/128 series
(49.22\%) satisfy this test; the equal-family proportion is 50.52\%.
At $H=1,000$, 94/128 (73.44\%) satisfy it, with equal-family proportion
72.87\%. All certified winners agree with the corresponding unique
sampled winner. No sampled winner set is certified strictly suboptimal
\emph{by these partial-grid intervals}. The intervals confirm a subset
of choices but leave the erroneous choices unresolved.

The strong baseline retains every distinct score class by requesting
$H=n$ and deduplicating complete threshold states. Its 1,536 calls also
complete, and certify the unique exhaustive winner on \emph{all 128}
series. Compared with the original $H=250$ selector, they prove ten
strictly suboptimal sampled winner sets and eight with regret at least
0.01. The largest guaranteed regret lower bound is 0.3985750301665.
At $H=1,000$, one sampled set is provably suboptimal, with guaranteed
regret 0.0018926253; none reaches 0.01. These are exact integer-endpoint
decisions, not comparisons to a floating-point reference with a tolerance.

\begin{table}[t]
\centering\small
\begin{tabular}{lrrrr}
\toprule
Retained grid & median states & certified winners & median width & max width\\
\midrule
250 ranks & 50 & 63/128 & $6.237\times10^{-3}$ & $\approx .99$\\
1,000 ranks & 68 & 94/128 & $3.900\times10^{-4}$ & $\approx .99$\\
All score classes & 96 & 128/128 & $2.766\times10^{-11}$ & $1.714\times10^{-9}$\\
\bottomrule
\end{tabular}
\caption{Integer enclosures at $L=30$. State counts and interval widths
are per score vector; winner guarantees compare all twelve models jointly.
A small median interval width does not replace the all-competitor test.}
\label{tab:certificates}
\end{table}

All prior certified winners are preserved, and no full-grid interval is
disjoint from its corresponding partial-grid interval. Independently
rounded intervals need not nest; real-arithmetic refinement monotonicity
is not asserted for the raw rounded outputs. Full-grid certification is
an explicit same-guarantee baseline, not an implicit claim that partial
certificates are more efficient. Costs are evaluated separately below.

\paragraph{A real target versus a floating-point reference.}
One $H=1,000$ interval has the floating-point exhaustive answer
$7.63\times10^{-16}$ below its lower endpoint. Direct analysis resolves
the case exactly: on the YAHOO/CNN input, three original intervals never
merge over widths zero through thirty, eight leading originals cover all
components, and the ninth original joins one negative in a complete score
tie. Buffer points enter only after recall saturates. At every width the
ideal value is $8/9+(1/9)(9/10)=89/90$, which the certificate contains.
The discrepancy is floating-point accumulation in the reference, not an
interval violation. No tolerance or widened interval was used to remove it.

% Sensitivity and timing tables are inserted after the frozen collections end.

\subsection{Sensitivity and matched computation costs}
All 36,864 declared numerical calls finish successfully, including all
4,608 upstream $H=250$ parity comparisons across the three widths.
Table~\ref{tab:budget-sensitivity} reports every declared budget/width
combination. Raising the budget to 1,000 removes practical selection
regret at widths 30 and 100 on this panel, while width 300 retains one
case with regret 0.095202. At $H=250$, widths 100 and 300 affect five
and three families respectively. Width 300's six cases fall below the
prespecified 5\% fraction; neither width replaces the primary result.

\begin{table}[t]
\centering\small
\begin{tabular}{lrrrrrr}
\toprule
Maximum width & $H=32$ & $64$ & $128$ & $250$ & $500$ & $1{,}000$\\
\midrule
30 (primary) & 21 & 16 & 11 & 8 & 1 & 0\\
100 & 20 & 15 & 12 & 7 & 2 & 0\\
300 & 20 & 15 & 11 & 6 & 2 & 1\\
\bottomrule
\end{tabular}
\caption{Series with best-case exhaustive regret at least 0.01; all
denominators are 128. These sensitivity cells are not independent studies.}
\label{tab:budget-sensitivity}
\end{table}

Timings use the fixed 30-series validation subset, all twelve models,
five independently randomized paired repetitions, one compute thread,
a 30-second per-call process watchdog and non-study warmup. Other
research compute jobs stop for this phase. Every timed scalar call includes
public-API validation, preparation and integration. Input decoding,
imports/JIT warmup, process setup, winner aggregation and record writing
are outside its timer. The broader collection wall time is recorded
separately and is not substituted for scalar-call latency. All 9,000
planned timing calls complete with output parity and no timeout.

Table~\ref{tab:matched-costs} reports per-case median times aggregated as
geometric comparator/exhaustive-radix ratios, with source-family bootstrap
intervals. Exact computation is modestly faster than matched compiled
sampling: the ratio is 1.044, with interval [1.029, 1.062]. The larger
ratio against unchanged Klein includes implementation differences and
cannot be interpreted as an intrinsic exact-versus-approximate gain.
The independent dense control additionally computes ROC internally;
its output scope is disclosed rather than used as a matched PR-only
speed claim. The saturation-exact interval includes one. These are warm,
one-host measurements on the released quantized score panel.

\begin{table}[t]
\centering\small
\begin{tabular}{lrrl}
\toprule
Comparator / exhaustive radix & geometric ratio & family 95\% interval & outputs\\
\midrule
Matched compiled, $H=250$ & 1.044 & [1.029, 1.062] & PR\\
Saturation exact & 1.011 & [0.999, 1.026] & PR\\
Unchanged Klein, $H=250$ & 5.077 & [4.186, 6.403] & PR\\
Independent dense exact & 13.481 & [12.505, 14.583] & ROC and PR\\
\bottomrule
\end{tabular}
\caption{360 paired series/model cases at $L=30$, five repetitions per
method. Ratios above one favor exhaustive radix. The matched compiled
row controls for preparation and implementation family.}
\label{tab:matched-costs}
\end{table}

% The separately frozen three-method certificate-cost collection follows.

\paragraph{Cost of the stronger guarantee.}
A separate frozen addendum measures the $H=250$ integer certificate,
the full distinct-score integer certificate, and floating-point exhaustive
radix together on the same 360 cases, with five randomized repetitions.
The three-method comparison was specified before any timing outcomes;
full-grid accuracy outcomes were already known when its final protocol
was frozen. The original two-method protocol remains unexecuted and
preserved, avoiding duplicate timing selection. All 5,400 calls pass,
with no missing records, failures, timeouts or output mismatch. Integer
endpoints match the completed certificate studies exactly.

\begin{table}[t]
\centering\small
\begin{tabular}{lrrr}
\toprule
Cost ratio & median & geometric mean & equal-family geometric mean\\
\midrule
250-rank certificate / full certificate & .826 & .815 & .820\\
250-rank certificate / numerical exact & 2.063 & 2.683 & 2.667\\
Full certificate / numerical exact & 2.713 & 3.290 & 3.254\\
\bottomrule
\end{tabular}
\caption{Separate three-way timing, same 360 score vectors and five
repetitions at $L=30$. These ratios describe time, so values below one
favor the numerator method. They are descriptive finite-panel summaries;
no confidence interval was specified for this addendum.}
\label{tab:certificate-cost}
\end{table}

The partial certificate saves approximately 18.5\% in the geometric
comparison with full certification, while certifying fewer winners on
the full panel. Both cost more than numerical exhaustive evaluation.
Summing the 360 per-case median times gives 4.686 s, 5.638 s and 3.127 s
for partial certificate, full certificate and numerical exact respectively;
these sums are not direct end-to-end collection timings. Actual collection
wall time is 71.918 s, including input/output, setup and orchestration.
The main and certificate timing runs use an Apple M1 Pro with 16 GiB RAM,
Python 3.12.4, NumPy 1.26.4 and Numba 0.59.1. A single active desktop host
limits portability; no cross-platform certificate performance claim is made.

The practical recommendation for these stored scores is therefore simple:
use the complete score classes for ordinary numerical evaluation, and
request full-grid integer enclosures when a rigorous ideal-target order
is required. Partial-grid certification provides a sufficient test and an
explicit inconclusive outcome, but this experiment does not establish a
better selection algorithm or a faster universal route to that guarantee.
